\documentclass[11pt]{article}
\usepackage{actstyle}
\settitle{ACT Reading: The Complete Study Guide}{Volume 1 --- Format, Method, and the Eleven Rules}

\begin{document}
\maketitleblock

\begin{abstract}
This guide covers everything the Enhanced ACT Reading section tests: 36 questions in 40 minutes across four passage units (literary narrative, social science, humanities, and natural science), including one paired set of shorter passages and passages that pair text with a chart or table. It is organized around ACT's three reporting categories---Key Ideas \& Details, Craft \& Structure, and Integration of Knowledge \& Ideas---which are broken into eleven rules. Each rule gives the question stems to recognize, a method, the traps to expect, and a drill set built on original passages. The guide ends with a wrong-answer field guide, a pacing protocol, a one-page quick reference, an error log, and a complete answer key. Volume~2 is a full-length, timed practice section built to the same specifications.
\end{abstract}

\scoreline{\textbf{Drills:} Rules 1--11}

\vspace{-4pt}
{\small\tableofcontents}

% =====================================================================
\section{Orientation: The Section at a Glance}
% =====================================================================

\subsection{Format}
\begin{center}
\begin{tabular}{@{}p{0.30\linewidth}p{0.62\linewidth}@{}}
\toprule
\textbf{Feature} & \textbf{Enhanced ACT Reading}\\
\midrule
Questions & 36 multiple-choice, 4 answer choices each\\
Scored questions & 27; the other 9 are unscored field-test items that look identical\\
Time & 40 minutes --- about 67 seconds per question\\
Passage units & 4 units, 9 questions each --- about 10 minutes per unit\\
Passage length & About 650--750 words; a paired set splits that length into two shorter passages\\
Guessing & No penalty --- never leave a question blank\\
Score & 1--36 scale; Reading is one-third of the composite (with English and Math)\\
Position & Third required section, after English and Math\\
\bottomrule
\end{tabular}
\end{center}

Because you cannot tell which questions are field-test items, \textbf{treat every question as scored.}

\subsection{The passage mix}
Every form includes one \textbf{literary narrative} (fiction or memoir) and \textbf{informational} passages drawn from the \textbf{social sciences, humanities, and natural sciences}. No outside knowledge is required---and outside knowledge is often a trap. Passages appear in three formats:

\begin{itemize}
  \item \textbf{Single passage} --- one excerpt, nine questions.
  \item \textbf{Paired passages} --- two shorter, related passages (Passage~A and Passage~B) with questions on each and on both.
  \item \textbf{Visual/quantitative (VQI) passage} --- text accompanied by a chart, graph, or table that you must interpret in the context of the passage.
\end{itemize}

You will see one of three combinations, and you will not know the order until the section begins:

\begin{center}
\begin{tabular}{@{}lccc@{}}
\toprule
\textbf{Combination} & \textbf{Single} & \textbf{Paired} & \textbf{VQI}\\
\midrule
Option 1 & 3 & 1 & ---\\
Option 2 & 3 & --- & 1\\
Option 3 & 2 & 1 & 1\\
\bottomrule
\end{tabular}
\end{center}

\subsection{Reporting categories}
Your score report breaks Reading into three categories. Shares are of the scored questions.

\begin{center}
\small
\begin{tabular}{@{}p{0.22\linewidth}ccp{0.44\linewidth}@{}}
\toprule
\textbf{Category} & \textbf{Share} & \textbf{Scored items} & \textbf{Skills tested (rules in this guide)}\\
\midrule
Key Ideas \& Details (KID) & 44--52\% & 12--14 & Close reading of stated details; inference; central ideas, themes, and summaries; relationships such as cause--effect, sequence, and comparison (Rules 1--4)\\
\addlinespace
Craft \& Structure (C\&S) & 26--33\% & 7--9 & Word meaning and word choice; text structure and the function of parts; author's purpose and point of view (Rules 5--8)\\
\addlinespace
Integration of Knowledge \& Ideas (IKI) & 19--26\% & 5--7 & Arguments---claims, evidence, reasoning; synthesis across paired passages; integrating text with charts and tables (Rules 9--11)\\
\bottomrule
\end{tabular}
\end{center}

\subsection{What changed from the legacy ACT}
\begin{center}
\small
\begin{tabular}{@{}lll@{}}
\toprule
& \textbf{Legacy ACT} & \textbf{Enhanced ACT}\\
\midrule
Questions / time & 40 in 35 min & 36 in 40 min\\
Time per question & $\approx$ 53 sec & $\approx$ 67 sec\\
Questions per passage & 10 & 9\\
Passage length & Longer & Shorter (about 650--750 words)\\
Key Ideas \& Details & 53--60\% & 44--52\%\\
Craft \& Structure & 25--30\% & 26--33\%\\
Integration of Knowledge \& Ideas & 15--23\% & 19--26\%\\
Charts/tables in Reading & No & Yes (VQI passages)\\
Unscored field-test items & Separate & Embedded in the section\\
\bottomrule
\end{tabular}
\end{center}

\textbf{What this means for you.} The test is less of a speed race and more of a precision test. More time per question means fewer careless misses are tolerated at the top end, and the larger Integration share means paired-passage, argument, and data questions now decide high scores.

% =====================================================================
\section{The Method}
% =====================================================================

\subsection{Reading the passage: map it, don't memorize it}
You cannot remember 700 words, and you don't need to. You need a \emph{map}: where things are, and what the passage is doing.

\begin{enumerate}
  \item \textbf{Read the introductory note.} The italic blurb tells you the source, genre, and sometimes the narrator. Ten seconds, every time.
  \item \textbf{Read the whole passage once, actively} (3--4 minutes). Do not stop to solve confusing sentences; mark them and move on.
  \item \textbf{Write a 3--6 word note beside each paragraph} (or highlight one phrase on the digital test): what the paragraph \emph{does}, not every detail it contains. Example: ``\emph{Ruskin: restoration = lie}.''
  \item \textbf{Mark the pivots:} contrast words (\emph{but, however, yet}), opinions (\emph{I believe, ought, unfortunately}), shifts in time, and any numbers.
  \item \textbf{State the Big Picture in one sentence} before you look at a question: the main idea plus the author's or narrator's attitude.
\end{enumerate}

\subsection{Answering: the four-step loop}
\begin{rulebox}{The Loop --- run it on every question}
\begin{enumerate}[leftmargin=1.6em]
  \item \textbf{Read the stem and classify it.} What rule is this? What is the locator---a line, a paragraph, a name, a keyword?
  \item \textbf{Go back to the text.} Read the referenced lines \emph{plus a sentence before and after}. For questions with no reference, use your paragraph notes.
  \item \textbf{Predict.} Put the answer in your own words before reading the choices. A vague prediction is fine; no prediction is not.
  \item \textbf{Eliminate.} Every wrong answer has at least one fatal word. Find it, cross it out, and choose the survivor that matches the text.
\end{enumerate}
\end{rulebox}

\textbf{The evidence standard.} If you cannot point to the words in the passage that make an answer right, it is not the answer. Correct answers are usually \emph{paraphrases} of the text; wrong answers often reuse the text's exact words in the wrong way.

\subsection{Order of attack within a passage}
\begin{enumerate}
  \item Questions with line or paragraph references (fast, contained).
  \item Detail and vocabulary questions.
  \item Inference, purpose, and structure questions.
  \item Big-picture and Integration questions last---by then you have reread most of the passage.
\end{enumerate}

\subsection{Signal words}
\begin{center}
\small
\begin{tabular}{@{}p{0.2\linewidth}p{0.45\linewidth}p{0.27\linewidth}@{}}
\toprule
\textbf{Signal} & \textbf{Words} & \textbf{What it tells you}\\
\midrule
Contrast & but, however, yet, although, still, nevertheless, instead, actually & The important idea usually comes \emph{after} the pivot\\
Concession & admittedly, granted, to be sure, of course \dots\ but & The author grants a point, then argues past it\\
Continuation & also, moreover, furthermore, likewise, similarly & Same direction; adds support\\
Cause/effect & because, so, therefore, thus, as a result, consequently & Map the arrow for Rule 4\\
Example & for instance, consider, in one study & Evidence for the claim just made (Rule 9)\\
Emphasis & indeed, in fact, crucially, above all & Likely a main point\\
\bottomrule
\end{tabular}
\end{center}

% =====================================================================
\section{Passage-Type Playbooks}
% =====================================================================

\subsection{Literary Narrative (fiction or memoir)}
\textbf{Read for:} who the characters are, what they want, how they feel about each other, and \emph{what changes}. Most literary passages turn on a shift---a realization, a reconciliation, a loss.\\
\textbf{Watch for:} a \emph{retrospective narrator} (an adult looking back on youth) whose present-day judgment differs from what they felt then; feelings shown through gestures rather than stated.\\
\textbf{Common questions:} character inference, narrator's point of view, the function of a scene or detail, the meaning of a figurative phrase.\\
\textbf{Trap:} choices that are too dramatic. ACT characters feel \emph{uneasy} or \emph{wistful} far more often than \emph{furious} or \emph{devastated}.

\subsection{Social Science}
\textbf{Read for:} the central claim or explanation, the studies and data used to support it, and \emph{whose} view each sentence represents (researcher, critic, author).\\
\textbf{Common questions:} cause--effect, details about studies, the author's claim versus others' claims, paired-passage synthesis.\\
\textbf{Trap:} attributing a quoted researcher's view to the author, or the reverse.

\subsection{Humanities}
\textbf{Read for:} the author's attitude. Humanities passages (art, architecture, literature, music, philosophy, often in the first person) are usually argumentative or reflective even when they sound descriptive.\\
\textbf{Common questions:} purpose, tone, point of view, word choice, the role of a personal anecdote.\\
\textbf{Trap:} ignoring a shift in the author's view (``I used to think \dots\ but'').

\subsection{Natural Science}
\textbf{Read for:} the phenomenon, the explanation, the competing hypotheses, and the evidence that decides between them. Terms are defined in the passage---use the passage's definition, not yours.\\
\textbf{Common questions:} detail, cause--effect chains, sequence of a process, function of an example, VQI data questions.\\
\textbf{Trap:} outside knowledge. If the passage doesn't say it, it isn't the answer---even if it's true.

\subsection{Paired Passages}
\begin{enumerate}
  \item Read Passage A; answer the Passage A questions.
  \item Read Passage B; answer the Passage B questions.
  \item Before the ``both passages'' questions, write the relationship in one sentence: \emph{agree? disagree? B gives an example of A? B concedes A's point but reaches a different conclusion?}
\end{enumerate}
This keeps you from mixing up which author said what---the single most common paired-passage error.

\subsection{Visual/Quantitative (VQI) Passages}
\begin{enumerate}
  \item Read the figure like a passage: \textbf{title, axes, units, legend, groups}. Note what is \emph{not} shown.
  \item For data-only questions, use only the figure. Check units (thousands? percent? kHz?).
  \item For integration questions, find the sentence in the text and the part of the figure that say the same thing.
  \item Distinguish \textbf{level} (how high) from \textbf{change} (how much it moved), and \textbf{correlation} from \textbf{cause}.
\end{enumerate}

% =====================================================================
\section{Key Ideas \& Details (44--52\%)}
% =====================================================================
The largest category. These questions ask what the passage says, what it implies, what it is mostly about, and how its ideas connect.

% ---------------------------------------------------------------------
\subsection{Rule 1 --- Stated details}
\begin{rulebox}{Rule 1: The answer is in the text, paraphrased}
Go back and find it---never answer a detail question from memory. The right choice restates the text; it does not add to it.
\end{rulebox}
\textbf{Stems:} \emph{According to the passage \dots; The passage states that \dots; The narrator indicates \dots; Which of the following is NOT mentioned \dots}\\
\textbf{Method:} pull a keyword from the stem, scan your paragraph notes, reread the sentence, predict, match.\\
\textbf{Traps:} the \emph{plausible but unstated} answer (true in real life, absent from the passage); the \emph{misplaced detail} (real words from the wrong paragraph). For EXCEPT/NOT questions, check off each choice you find in the text; the leftover is the answer.

\passagehead{DRILL 1}{SOCIAL SCIENCE}{This passage is original and describes a fictional town.}
\begin{passage}
When the Marlow Falls Public Library added a shelf of lendable tools in 2019, the head librarian, Dana Ruiz, expected the drills and ladders to be the most popular items. They were not. In the program's first year, the most frequently borrowed item was a soil-testing kit, followed by a sewing machine and a set of folding tables. Ruiz attributes the pattern to cost and frequency: people who drill holes do so often enough to buy a drill, she argues, while a gardener might test soil once a season and a family might need twelve folding tables once a year. The library now chooses new items by surveying patrons about what they would use ``rarely but urgently.'' Late fees, which the library eliminated for books in 2021, still apply to tools, a policy Ruiz defends on the grounds that a missing tile saw, unlike a missing novel, can delay another family's repair.
\end{passage}

\Q{KID}{C}{R1 Detail}{``the most frequently borrowed item was a soil-testing kit.'' Drills (A) were what Ruiz \emph{expected}; B and D ranked lower.}
{According to the passage, the most frequently borrowed item in the program's first year was:}
{a drill.}{a sewing machine.}{a soil-testing kit.}{a set of folding tables.}

\Q{KID}{A}{R1 Detail}{``Ruiz attributes the pattern to cost and frequency.'' B--D are never mentioned.}
{The passage states that Ruiz explains the borrowing pattern in terms of:}
{how much items cost and how often people need them.}{how heavily the library advertised each item.}{how heavy or difficult to carry each item is.}{the ages of the patrons who borrowed most often.}

\Q{KID}{B}{R1 Detail; plausible-but-unstated trap}{The stated reason: a missing tool ``can delay another family's repair.'' A is plausible in real life but is never stated.}
{According to the passage, the library still charges late fees for tools because:}
{tools are more expensive to replace than books.}{a late tool can hold up another patron's project.}{the fees pay for new items on the tool shelf.}{patrons asked the library to keep the fees.}

% ---------------------------------------------------------------------
\subsection{Rule 2 --- Inference}
\begin{rulebox}{Rule 2: An inference is the smallest step the text forces you to take}
Choose the answer that \emph{must} be true given the text, not the one that \emph{could} be true. If you need a story to connect the text to the choice, the choice is wrong.
\end{rulebox}
\textbf{Stems:} \emph{It can reasonably be inferred \dots; The passage most strongly suggests \dots; Based on the passage, X most likely \dots}\\
\textbf{Method:} find the lines; ask ``what does this \emph{show} without saying?'' Predict modestly.\\
\textbf{Traps:} over-reaching (a big conclusion from a small clue); the literal restatement that doesn't answer what is implied; emotional exaggeration.

\passagehead{DRILL 2}{LITERARY NARRATIVE}{This passage is an original story.}
\begin{passage}
Ellis had rehearsed the speech on the bus, twice, and once more in the elevator, so that by the time he reached the fourth floor the words had worn smooth as coins. He knocked. His sister opened the door holding a dish towel, and for a moment neither of them said anything at all. Behind her, the apartment looked exactly as he had imagined it would---the same faded rug their mother had refused to throw out, the same crooked lamp---except that the walls, which in his imagination had been bare, were covered with his own drawings from childhood, framed.

``You're early,'' she said. He had meant to begin with an apology. Instead he heard himself ask where she had found the one of the lighthouse, and she laughed, and stepped aside to let him in.
\end{passage}

\Q{KID}{A}{R2 Inference}{A rehearsed speech, a planned apology, and a long silence at the door all point to time apart and something unresolved. B is contradicted (he visits); C is contradicted (they share a mother); D over-reads one detail.}
{It can most reasonably be inferred that Ellis and his sister:}
{had been out of contact, with some unresolved tension between them.}{had lived together in the apartment until recently.}{were meeting for the first time as adults.}{disagreed about whether to keep their mother's rug.}

\Q{KID}{D}{R2 Inference}{Framing and hanging his childhood drawings shows she kept valuing him. Nothing supports selling (A), resentment (B), or a hobby (C).}
{The detail that the sister's walls were covered with Ellis's framed childhood drawings most strongly suggests that she:}
{was preparing to sell the drawings.}{resented Ellis for leaving the drawings behind.}{had recently taken up framing as a hobby.}{had continued to value her connection to Ellis.}

\Q{KID}{B}{R2 Inference}{He abandons the rehearsed apology for a spontaneous question; she laughs and lets him in---the meeting turns warmer and less scripted. A over-reaches: he may still feel sorry.}
{Ellis's question about the lighthouse drawing, in place of the apology he had planned, suggests that:}
{he no longer believed he owed his sister an apology.}{the meeting took a warmer, less rehearsed turn than he expected.}{he had forgotten his speech because he was nervous.}{his sister had asked him not to apologize.}

% ---------------------------------------------------------------------
\subsection{Rule 3 --- Central ideas, themes, and summaries}
\begin{rulebox}{Rule 3: The main idea covers the whole---no less, no more}
A main-idea answer must account for every paragraph (or the whole paragraph asked about). Reject choices that are too narrow (one detail), too broad (beyond the passage), or distorted (right topic, wrong claim).
\end{rulebox}
\textbf{Stems:} \emph{The main idea of the passage is \dots; The passage is primarily concerned with \dots; The main point of the third paragraph is \dots; Which best summarizes \dots}\\
\textbf{Method:} use your Big Picture sentence. Test each choice: ``Does this describe paragraph 1? The last paragraph?''\\
\textbf{Traps:} a true detail posing as the main idea; a choice with the right topic but an absolute word (\emph{always, stopped, all}).

\passagehead{DRILL 3}{HUMANITIES}{This passage is original.}
\begin{passage}
For most of the twentieth century, museums treated the backs of paintings as storage space---a place for inventory numbers, shipping labels, and the occasional pencil note from a framer. Conservators now treat them as documents. A customs stamp can reveal when a painting crossed a border; a dealer's label can fill a gap in its ownership history; a sketch on the reverse of a canvas can show what an artist abandoned before beginning the work on the front.

This shift has practical consequences. When a painting is relined---backed with a new canvas to strengthen it---the old reverse is often covered permanently. Some museums now photograph every reverse before such treatment, and a few have begun displaying selected paintings in freestanding cases so that visitors can walk around them. The back of a painting, once the least looked-at part of the object, is becoming part of what the museum shows.
\end{passage}

\Q{KID}{C}{R3 Central idea}{Paragraph 1 = backs are now treated as historical documents; paragraph 2 = this changes treatment and display. Only C covers both. B and D are details; A is never claimed.}
{Which of the following best states the main idea of the passage?}
{Relining damages paintings and should be abandoned by museums.}{Customs stamps are the most reliable source of a painting's ownership history.}{Museums have come to value the backs of paintings as historical records, which has changed how they treat and display them.}{Artists frequently abandoned sketches on the backs of their canvases.}

\Q{KID}{D}{R3 Paragraph main idea}{Paragraph 2 opens ``This shift has practical consequences'' and lists them. A, B, and C each distort one detail.}
{The main idea of the second paragraph is that:}
{relining is the best way to strengthen an old painting.}{visitors prefer freestanding cases to traditional wall displays.}{photographs of a painting's reverse can replace the original.}{the new view of paintings' backs has led to changes in museum practice.}

\Q{KID}{A}{R3 Summary}{A captures both paragraphs accurately. B: ``stopped'' is too strong. C: wrong time period. D: visitor preference is never measured.}
{Which summary of the passage is most accurate?}
{Once overlooked, the backs of paintings are now studied for the history they record, and some museums preserve and even exhibit them.}{Museums have stopped relining paintings because the process covers important labels.}{Conservators first learned to read customs stamps in the early twentieth century.}{Most museum visitors now prefer to see the backs of paintings.}

% ---------------------------------------------------------------------
\subsection{Rule 4 --- Relationships: cause, sequence, comparison}
\begin{rulebox}{Rule 4: Map the arrow}
For every relationship question, write the arrow: \emph{A causes B}, \emph{A happens before B}, \emph{A is larger than B}. Wrong answers reverse the arrow, swap in a different cause, or invent a link the text never makes.
\end{rulebox}
\textbf{Stems:} \emph{According to the passage, X occurs because \dots; As a result of X \dots; Which happened first \dots; Compared with X, Y \dots}\\
\textbf{Method:} find the signal word (\emph{because, so, causes, as a result, before, by contrast}). Draw the chain in the margin.\\
\textbf{Traps:} reversed causation; correlation sold as cause; a real cause from the passage attached to the wrong effect.

\passagehead{DRILL 4}{NATURAL SCIENCE}{This passage is original.}
\begin{passage}
Each spring, the snowpack on a mountain range acts as a reservoir, releasing water gradually as it melts. In recent decades, warmer winters in many western ranges have caused more precipitation to fall as rain rather than snow. Rain runs off immediately, so rivers fed by these ranges now peak earlier in the year---sometimes in March rather than June---and carry less water in late summer, when farms and cities need it most.

Dust compounds the problem. When wind deposits dust from dry lowlands onto the snow, the darkened surface absorbs more sunlight and melts sooner, in some basins by several weeks. Dust can therefore shift the timing of runoff even in a year when total snowfall is normal. By contrast, a cool spring with heavy late-season storms can delay melt and partly offset both effects, though such springs have become less common.
\end{passage}

\Q{KID}{B}{R4 Cause--effect}{Chain: warmer winters $\rightarrow$ more rain $\rightarrow$ immediate runoff $\rightarrow$ earlier peak. C reverses the text (such springs are \emph{less} common); D: dust does not reduce snowfall.}
{According to the passage, rivers fed by these mountain ranges now peak earlier in the year because:}
{farms and cities withdraw more water in the spring.}{more precipitation falls as rain, which runs off immediately.}{late-season storms have become more common.}{dust reduces the total amount of snowfall.}

\Q{KID}{D}{R4 Cause--effect}{``the darkened surface absorbs more sunlight and melts sooner.'' The others invent mechanisms.}
{According to the passage, dust on snow speeds melting because it:}
{adds weight that compresses the snowpack.}{reduces the amount of snow that falls.}{carries warm air from the lowlands.}{darkens the surface so that it absorbs more sunlight.}

\Q{KID}{A}{R4 Comparison}{``By contrast'' signals the opposite effect: cool, stormy springs \emph{delay} melt while dust \emph{advances} it. C reverses ``less common.''}
{The passage indicates that, compared with dust deposits, a cool spring with heavy late-season storms:}
{has the opposite effect on the timing of snowmelt.}{has a larger effect on total snowfall.}{occurs more often than it used to.}{affects only rivers in dry lowlands.}

\Q{KID}{C}{R4 Sequence}{Paragraph 1 gives the order exactly. A reverses dust's effect; B and D put effects before causes.}
{Which sequence of events is consistent with the passage?}
{Dust falls $\rightarrow$ snowfall increases $\rightarrow$ rivers peak later}{Late-summer water shortages $\rightarrow$ warmer winters $\rightarrow$ more rain}{Warmer winters $\rightarrow$ more rain and less snow $\rightarrow$ earlier river peaks $\rightarrow$ less late-summer flow}{Earlier river peaks $\rightarrow$ more snow $\rightarrow$ cooler springs}

% =====================================================================
\section{Craft \& Structure (26--33\%)}
% =====================================================================
These questions ask \emph{how} the passage works: what words mean in context, why the author chose them, what a part of the passage is doing, and from whose point of view.

% ---------------------------------------------------------------------
\subsection{Rule 5 --- Word meaning in context}
\begin{rulebox}{Rule 5: Cover the word, supply your own, then match}
Read the sentence with the word blanked out and fill in a simple word that fits. The most common dictionary meaning of the tested word is usually the trap.
\end{rulebox}
\textbf{Stems:} \emph{As it is used in line X, the word ``\_\_'' most nearly means \dots}\\
\textbf{Method:} read the full sentence and the one before it; predict; plug each choice back in.\\
\textbf{Traps:} the literal or most familiar meaning; a choice that fits the word but not the sentence's tone.

\passagehead{DRILL 5}{SOCIAL SCIENCE}{This passage is original and describes a fictional town.}
\begin{passage}
The town council's plan to reopen the old rail depot as a market drew a \linelabel{d5cool}cool reception at first. Residents worried about parking, about noise, about the cost. But the council's chair, who had spent thirty years running a hardware store on Main Street, understood that objections are often a form of attention, and attention can be \linelabel{d5cult}cultivated. She held six open meetings in the depot itself, let residents walk the platform, and posted every estimate---including the unflattering ones---on the town website. By the time the council voted, the plan's most vocal critics had become its most \linelabel{d5exact}exacting supporters, pressing for details on vendor fees and weekend hours rather than whether the market should exist at all.
\end{passage}

\Q{CS}{B}{R5 Word in context; literal trap}{Residents ``worried'' about parking, noise, cost---the reception was unenthusiastic. C is the literal (temperature) trap; A is the slang trap.}
{As it is used in \Line{d5cool}, the word \emph{cool} most nearly means:}
{fashionable.}{unenthusiastic.}{slightly cold.}{composed.}

\Q{CS}{C}{R5 Word in context}{The chair grows residents' attention through meetings and openness---\emph{developed}. A is the farming trap; B (``cultivated taste'') is the other dictionary trap.}
{As it is used in \Line{d5cult}, the word \emph{cultivated} most nearly means:}
{farmed.}{refined.}{developed.}{tolerated.}

\Q{CS}{D}{R5 Word in context}{Supporters ``pressing for details on vendor fees and weekend hours'' are demanding and detail-focused. B plays on ``exact revenge.''}
{As it is used in \Line{d5exact}, the word \emph{exacting} most nearly means:}
{punishing.}{vengeful.}{reluctant.}{demanding and attentive to detail.}

% ---------------------------------------------------------------------
\subsection{Rule 6 --- Word choice, tone, and connotation}
\begin{rulebox}{Rule 6: Ask what the word does that a neutral synonym wouldn't}
Authors choose words for their associations. Identify the connotation, then pick the effect it creates. Tone answers are usually moderate---two words like \emph{wry} or \emph{rueful}, not \emph{furious}.
\end{rulebox}
\textbf{Stems:} \emph{The author uses the word/phrase X most likely to \dots; The tone of the passage is best described as \dots; The quotation marks around X serve to \dots}\\
\textbf{Method:} swap in a plain synonym and ask what is lost.

\begin{center}
\small
\begin{tabular}{@{}p{0.29\linewidth}p{0.29\linewidth}p{0.34\linewidth}@{}}
\toprule
\textbf{Positive} & \textbf{Negative} & \textbf{Mixed / restrained}\\
\midrule
admiring, reverent, earnest, affectionate, enthusiastic, appreciative & critical, skeptical, sardonic, indignant, dismissive, disdainful & ambivalent, measured, wistful, nostalgic, rueful, wry, detached, objective\\
\bottomrule
\end{tabular}
\end{center}

\passagehead{DRILL 6}{HUMANITIES (MEMOIR)}{This passage is original.}
\begin{passage}
My grandfather's workshop was not so much organized as \linelabel{d6neg}negotiated. Every jar of screws, every coil of wire, had arrived at its place through some long-forgotten compromise between convenience and memory, and he could find anything in it in the dark. After he died, my aunt hired a company to \linelabel{d6clear}``clear the space.'' They were efficient. In an afternoon they reduced fifty years of negotiation to a swept concrete floor and four tidy contractor bags. I don't blame them; it was the job she asked for. \linelabel{d6last}But I have never since heard the word \emph{efficient} without hearing, beneath it, the faint scrape of a broom.
\end{passage}

\Q{CS}{A}{R6 Word choice}{The next sentence explains: each item's place came from ``compromise between convenience and memory''---many small decisions over time. B takes ``negotiated'' literally.}
{The author describes the workshop as ``negotiated'' (\Line{d6neg}) primarily to suggest that its arrangement:}
{developed gradually through the grandfather's many small decisions.}{was the subject of arguments between family members.}{was too disorganized to be useful.}{was planned by a professional organizer.}

\Q{CS}{C}{R6 Word choice; quotation marks}{Quotation marks mark the phrase as someone else's language and distance the author from its impersonal tidiness. The author says ``I don't blame them,'' so B is too harsh.}
{The author places ``clear the space'' (\Line{d6clear}) in quotation marks most likely to:}
{indicate that the phrase is a legal term.}{accuse the aunt of being dishonest.}{signal that the words belong to someone else and convey the author's distance from their impersonal tone.}{emphasize the author's full agreement with the decision.}

\Q{CS}{D}{R6 Tone}{Quiet regret without blame (``I don't blame them'')---rueful. A is too strong; B misses the loss; C misses the feeling.}
{The tone of the final sentence (\Line{d6last}) is best described as:}
{bitterly accusatory.}{cheerfully nostalgic.}{detached and clinical.}{quietly rueful.}

% ---------------------------------------------------------------------
\subsection{Rule 7 --- Text structure and the function of parts}
\begin{rulebox}{Rule 7: Function questions ask ``why is this here?''}
Name the job a sentence or paragraph does relative to its neighbors: introduces a claim, gives an example, offers a counterargument, shifts focus, concludes. Choices that merely restate content are wrong.
\end{rulebox}
\textbf{Stems:} \emph{The main function of the second paragraph is to \dots; The author includes X primarily to \dots; Which best describes the organization of the passage \dots}\\
\textbf{Function verbs to know:} \emph{introduce, illustrate, support, qualify, contrast, concede, refute, transition, summarize, complicate.}

\passagehead{DRILL 7}{SOCIAL SCIENCE}{This passage is original; the experiment described is a well-known study in behavioral economics.}
\begin{passage}
Economists long assumed that people value a dollar saved the same as a dollar earned, no matter the circumstances. Shoppers suggest otherwise.

In one well-known experiment, people were asked whether they would drive across town to save \$5 on a \$15 calculator. Most said yes. Asked whether they would make the same drive to save \$5 on a \$125 jacket, most said no. \linelabel{d7same}The savings were identical; only the comparison point had changed.

Findings like this helped launch behavioral economics, a field that treats such inconsistencies not as random errors but as predictable patterns. Its practical influence is now wide: retirement plans that enroll workers automatically, cafeterias that place fruit at eye level, and utility bills that compare a household's energy use with its neighbors' all draw on the idea that the way a choice is framed shapes the choice people make.
\end{passage}

\Q{CS}{B}{R7 Function of a paragraph}{Paragraph 1 claims shoppers contradict the old assumption; paragraph 2 is the example that shows it. It supports, not counters, the passage's claim (A).}
{The second paragraph functions mainly to:}
{introduce a counterargument to the passage's main claim.}{provide an example that illustrates the claim made in the first paragraph.}{explain the history of behavioral economics.}{show that shoppers' decisions are random.}

\Q{CS}{A}{R7 Function of a sentence}{The sentence points out what makes the two answers inconsistent: same \$5, different reference price.}
{The sentence in \Line{d7same} serves mainly to:}
{clarify what makes the shoppers' responses inconsistent.}{suggest that the experiment was poorly designed.}{introduce the field of behavioral economics.}{argue that the jacket was overpriced.}

\Q{CS}{C}{R7 Overall structure}{Assumption stated (\P1) $\rightarrow$ challenged with evidence (\P2) $\rightarrow$ linked to a field and its applications (\P3).}
{Which of the following best describes the overall structure of the passage?}
{Two theories are compared, and one is rejected.}{A personal anecdote is followed by a general rule.}{A traditional assumption is stated, challenged with evidence, and connected to a broader field and its applications.}{A problem is described, and several solutions are ranked.}

% ---------------------------------------------------------------------
\subsection{Rule 8 --- Purpose and point of view}
\begin{rulebox}{Rule 8: Who is speaking, when, and how do they feel?}
Identify the narrator or author, their vantage point (participant, observer, adult looking back), and their attitude. ACT authors are rarely neutral and rarely extreme: \emph{partly accepting but critical} beats \emph{hostile}.
\end{rulebox}
\textbf{Stems:} \emph{The author's main purpose is to \dots; The point of view from which the passage is told is that of \dots; The author's attitude toward X is best described as \dots; The author mentions X primarily to \dots}\\
\textbf{Method:} look for opinion language and concessions (\emph{not entirely wrong, but \dots}). Purpose answers begin with verbs: \emph{argue, explain, describe, challenge, explore.}

\passagehead{DRILL 8}{NATURAL SCIENCE}{This passage is original.}
\begin{passage}
It has become fashionable to describe the octopus as an alien intelligence---a mind that evolved along a path so separate from ours that understanding it may be impossible. The description is vivid, and it is not entirely wrong; the last common ancestor of octopuses and humans was probably a simple wormlike creature that lived more than 500 million years ago. But the ``alien'' label can excuse a kind of laziness. Octopuses solve problems, recognize individual people, and appear to play, and each of these behaviors can be studied with the same patient methods researchers use on \linelabel{d8crow}crows and chimpanzees. Calling a mind alien is a way of declaring the investigation finished before it has properly begun.
\end{passage}

\Q{CS}{D}{R8 Purpose}{The pivot ``But the `alien' label can excuse a kind of laziness'' and the final sentence state the purpose. C reverses it; A and B are never the goal.}
{The author's main purpose in the passage is to:}
{prove that octopuses are more intelligent than crows.}{explain the evolutionary history of octopuses.}{defend the description of octopuses as alien intelligences.}{argue that calling octopus minds ``alien'' discourages the careful study they deserve.}

\Q{CS}{B}{R8 Attitude}{``Vivid'' and ``not entirely wrong'' (partial acceptance) then ``laziness'' (criticism). D is too extreme; A and C ignore half the passage.}
{The author's attitude toward the ``alien intelligence'' description is best described as:}
{wholly approving.}{partly accepting but ultimately critical.}{indifferent.}{openly contemptuous.}

\Q{CS}{A}{R8 Purpose of a detail}{Crows and chimpanzees are animals already studied with ``patient methods''---the same methods could apply to octopuses.}
{The author mentions crows and chimpanzees (\Line{d8crow}) primarily to:}
{suggest that octopus behavior can be studied with established methods.}{argue that octopuses are closely related to birds and primates.}{show that octopuses are less intelligent than other animals.}{introduce a new topic for the rest of the passage.}

% =====================================================================
\section{Integration of Knowledge \& Ideas (19--26\%)}
% =====================================================================
The category the Enhanced ACT expanded. These questions ask you to evaluate arguments, connect two passages, and combine text with data.

% ---------------------------------------------------------------------
\subsection{Rule 9 --- Arguments: claims, evidence, reasoning}
\begin{rulebox}{Rule 9: Separate the claim from its support}
Find the author's claim first. Evidence answers must connect \emph{directly} to that specific claim. For strengthen/weaken questions, attack or reinforce the link between the evidence and the conclusion.
\end{rulebox}
\textbf{Stems:} \emph{The author's main claim is \dots; Which evidence most directly supports \dots; The statement X functions in the argument as \dots; Which of the following, if true, would most weaken \dots}\\
\textbf{Claim vs.\ evidence:} claims are arguable (\emph{should, may, better}); evidence is reported (studies, numbers, examples, quotations).\\
\textbf{Weakeners to look for:} a \emph{confound} (another cause), an unrepresentative sample, a measurement problem.

\passagehead{DRILL 9}{SOCIAL SCIENCE}{This passage is original and describes a fictional city.}
\begin{passage}
Cities that want more people to bicycle usually start by painting lanes. \linelabel{d9paint}Paint is cheap, fast, and visible, and it lets officials point to new miles of ``bike infrastructure'' at the end of a budget year. But paint does not change who rides. Surveys in many cities have found that most residents say they would cycle more if they felt safe, and that a painted line beside fast traffic does not make them feel safe. When the city of Lindgren replaced four miles of painted lanes with lanes separated from traffic by curbs, ridership on those streets roughly tripled within two years, while ridership on comparable painted-lane streets barely changed. Separated lanes cost more per mile. A city that builds fewer miles of them, however, may get more riders for its money than a city that paints many miles no one uses.
\end{passage}

\Q{IKI}{C}{R9 Identify the claim}{The final sentence states the claim: fewer miles of separated lanes may yield more riders per dollar. A and D are too extreme; B is never argued.}
{The author's main claim is that:}
{painted lanes are dangerous and should be removed.}{surveys of residents are unreliable.}{separated lanes may increase ridership more cost-effectively than painted lanes.}{cities should stop funding bicycle infrastructure.}

\Q{IKI}{B}{R9 Evidence}{The Lindgren comparison (tripled vs.\ barely changed) is the direct evidence that separated lanes get more riders. D actually cuts against the claim; A and C explain paint's appeal.}
{Which evidence does the author use most directly to support the main claim?}
{Paint is cheap and fast to apply.}{Ridership on Lindgren's separated-lane streets rose far more than on comparable painted-lane streets.}{Officials can point to new miles of lanes at the end of a budget year.}{Separated lanes cost more per mile.}

\Q{IKI}{D}{R9 Role of a statement}{The author concedes why paint is attractive before arguing against relying on it (``But paint does not change who rides'').}
{The statement in \Line{d9paint} (``Paint is cheap, fast, and visible'') functions in the argument as:}
{the author's main recommendation.}{evidence that painted lanes increase ridership.}{a statistic from a resident survey.}{an acknowledgment of why the approach the author criticizes is appealing.}

\Q{IKI}{A}{R9 Weaken}{A new campus on the same streets offers another cause for the ridership jump---a confound. B and C strengthen; D is irrelevant.}
{Which of the following, if true, would most weaken the author's argument?}
{During the same two years, a new university campus opened along the separated-lane streets and attracted thousands of students.}{Parents in Lindgren say they are more willing to let children ride on separated lanes.}{Surveys in other cities also found that safety concerns keep residents from cycling.}{Lindgren installed its painted lanes before it installed its separated lanes.}

% ---------------------------------------------------------------------
\subsection{Rule 10 --- Synthesis across paired passages}
\begin{rulebox}{Rule 10: Solve each passage alone, then name the relationship}
Answer each passage's own questions first. Before ``both passages'' questions, write one sentence describing how the passages relate. Then check every choice against \emph{both} texts.
\end{rulebox}
\textbf{Stems:} \emph{Both authors would most likely agree \dots; How would the author of Passage B respond to \dots; Which best describes the relationship between the passages \dots; Unlike Passage A, Passage B \dots}\\
\textbf{Traps:} a choice true of only one passage; a reversed attribution; ``agree'' choices that one author would merely tolerate.

\passagehead{DRILL 10}{PAIRED PASSAGES}{Both passages are original.}
\subpassage{Passage A}
\begin{passage}
The research on homework in the early grades is remarkably consistent: for young children, the link between time spent on homework and achievement is weak, and sometimes nonexistent. What homework reliably produces at that age is conflict---over the kitchen table, at the end of a long day, between tired children and tired parents. An eight-year-old's evening has better uses: reading for pleasure, playing outside, sleeping. Schools that have dropped homework in the elementary grades report few complaints once families adjust. The burden of proof belongs to those who would keep it.
\end{passage}
\subpassage{Passage B}
\begin{passage}[cont]
Critics of elementary homework are right about one thing: piles of worksheets do young children no good. But the case against homework usually measures it only by test scores, and test scores were never its main purpose in the early grades. Ten or fifteen minutes of practice teaches a child to start a task without being told and to finish it without being watched---habits that matter enormously later on. It also gives parents a nightly window into what their child is learning. The answer to bad homework is better homework, and less of it, not none at all.
\end{passage}

\Q{IKI}{C}{R10 Response across passages}{B says the anti-homework case ``measures it only by test scores, and test scores were never its main purpose.'' B never disputes the research (A).}
{The author of Passage B would most likely respond to Passage A's claim that homework has little link to achievement in the early grades by:}
{disputing the research Passage A describes.}{agreeing that homework should be eliminated.}{arguing that achievement is not the main purpose of homework at that age.}{claiming that test scores rise where homework is assigned.}

\Q{IKI}{A}{R10 Agreement}{A opposes homework; B concedes ``piles of worksheets do young children no good.'' Both reject heavy loads. B and C belong to one author only.}
{Both authors would most likely agree that:}
{large amounts of homework do not benefit young children.}{homework teaches young children independence.}{schools should eliminate homework in the elementary grades.}{parents should help with homework every night.}

\Q{IKI}{B}{R10 Relationship}{B concedes A's point (``right about one thing'') but concludes ``better homework \dots not none at all.'' C reverses the order---A never responds to B.}
{Which choice best describes the relationship between the two passages?}
{Passage B provides data that support Passage A's conclusion.}{Passage B concedes a point made in Passage A but argues for a different conclusion.}{Passage A responds directly to arguments made in Passage B.}{Both passages summarize research without taking a position.}

\Q{IKI}{D}{R10 Contrast}{A weighs achievement and family time (conflict, evenings); B weighs habits and parents' view into learning.}
{Passage A evaluates homework mainly by its effects on \_\_\_\_, whereas Passage B emphasizes its effects on \_\_\_\_.}
{cost; student enjoyment}{teacher workload; test scores}{long-term careers; class size}{achievement and family time; work habits and parents' insight}

% ---------------------------------------------------------------------
\subsection{Rule 11 --- Integrating text with charts and tables (VQI)}
\begin{rulebox}{Rule 11: Read the figure like a passage}
Title, axes, units, groups---before any question. Data-only questions need only the data. Integration questions need both the text and the figure, and nothing outside them.
\end{rulebox}
\textbf{Stems:} \emph{According to the table \dots; The data in the figure best support which claim from the passage \dots; Which additional information would best help evaluate \dots}\\
\textbf{Traps:} confusing level with change; ignoring units (thousands, percent); assuming the most recent year is the highest; choosing a claim the figure is \emph{consistent} with but does not \emph{show}.

\passagehead{DRILL 11}{SOCIAL SCIENCE (WITH TABLE)}{This passage and its data are original and describe a fictional library system.}
\begin{passage}
Libraries are often described as losing ground to screens. The Halvorsen County Library's figures complicate that story. In-person visits fell sharply in 2022, when branches operated on reduced hours, and have since recovered most---but not all---of the lost ground. Print checkouts have declined steadily. Yet e-book loans have grown so quickly that total circulation, print and digital combined, is now higher than it was a decade ago. The library's director argues that the building's role is shifting from warehouse to meeting place: fewer people come in to borrow a book, but more come in to use a study room, attend a class, or get help with a form.
\end{passage}

\begin{center}
\small
\textbf{Table 1.} Halvorsen County Library activity (in thousands)\\[3pt]
\begin{tabular}{@{}lccc@{}}
\toprule
\textbf{Year} & \textbf{In-person visits} & \textbf{Print checkouts} & \textbf{E-book loans}\\
\midrule
2016 & 410 & 520 & \phantom{0}60\\
2019 & 395 & 480 & 120\\
2022 & 250 & 400 & 230\\
2025 & 370 & 360 & 260\\
\bottomrule
\end{tabular}
\end{center}

\Q{IKI}{B}{R11 Data reading}{$260 \div 60 \approx 4.3$, so about 4 times. D confuses ``times'' with the difference in thousands (200).}
{According to Table 1, e-book loans in 2025 were approximately how many times their 2016 level?}
{2}{4}{6}{200}

\Q{IKI}{C}{R11 Data reading; most-recent-year trap}{Combined print + e-book: 2016 = 580; 2019 = 600; 2022 = 630; 2025 = 620. The peak is 2022, not the latest year.}
{Based on Table 1, in which year was combined print and e-book circulation highest?}
{2016}{2019}{2022}{2025}

\Q{IKI}{A}{R11 Text--data integration}{Visits: 395 $\rightarrow$ 250 $\rightarrow$ 370, below the 2019 level---``most but not all.'' B and C are interpretations the table does not measure; D describes others' view.}
{The data in Table 1 most directly support which statement from the passage?}
{``In-person visits \dots have since recovered most---but not all---of the lost ground.''}{``more come in to use a study room, attend a class, or get help with a form.''}{``the building's role is shifting from warehouse to meeting place.''}{``Libraries are often described as losing ground to screens.''}

\Q{IKI}{D}{R11 Evaluating a claim with data}{The director's claim is about meeting-place uses; only class attendance and room bookings measure that. A--C don't bear on how the building is used.}
{Which additional data would best help evaluate the director's argument that the library building is becoming a meeting place?}
{The number of e-book titles in the catalog each year}{The county's population in 2016}{The price the library pays for e-book licenses}{Yearly class attendance and study-room bookings}

% =====================================================================
\section{The Wrong-Answer Field Guide}
% =====================================================================
Wrong answers are written to a pattern. Learn the patterns and you can eliminate quickly and confidently.

\begin{center}
\small
\begin{tabular}{@{}p{0.2\linewidth}p{0.42\linewidth}p{0.3\linewidth}@{}}
\toprule
\textbf{Trap} & \textbf{What it looks like} & \textbf{The tell}\\
\midrule
Extreme & Right idea, too strong: \emph{always, never, proves, stopped, all} & Can you find language in the text that strong?\\
Out of scope & Plausible, even true, but not in the passage & ``Where does it say that?''---nowhere\\
Half-right & First half matches; second half adds an error & Read every word of the choice\\
Misplaced detail & Real words from the wrong paragraph or wrong person & Check the location in the stem\\
Reversal & Cause and effect, before and after, or more and less are swapped & Draw the arrow (Rule 4)\\
True but irrelevant & Accurate statement that doesn't answer this question & Reread the stem\\
Literal meaning & The dictionary sense of a word used figuratively & Cover the word and predict (Rule 5)\\
Too narrow / too broad & One detail, or more than the passage covers, as the main idea & Does it cover every paragraph? (Rule 3)\\
Wrong speaker & A cited expert's view attributed to the author, or Passage A's view to B & Who said it?\\
Consistent, not shown & A claim the data don't contradict but don't demonstrate & Point to the cell or bar that proves it\\
\bottomrule
\end{tabular}
\end{center}

\textbf{The 2-to-1 rule.} When you are stuck between two choices, find the single word that differs and go back to the text for \emph{that word only}. Do not reread the whole paragraph.

% =====================================================================
\section{Pacing and Test-Day Protocol}
% =====================================================================

\subsection{The 10-minute unit}
\begin{center}
\small
\begin{tabular}{@{}lcl@{}}
\toprule
\textbf{Phase} & \textbf{Time} & \textbf{Goal}\\
\midrule
Intro note + active read + paragraph notes & 3--4 min & Map and Big Picture sentence\\
Line-reference and detail questions & 3 min & Fast, evidence-based points\\
Inference, function, Integration questions & 3 min & Careful elimination\\
Buffer / flag review & 0--1 min & Return to one flagged question\\
\bottomrule
\end{tabular}
\end{center}

\subsection{Clock checkpoints}
\begin{center}
\small
\begin{tabular}{@{}lcccc@{}}
\toprule
\textbf{After unit} & 1 & 2 & 3 & 4\\
\midrule
Time used (target) & 10:00 & 20:00 & 30:00 & 40:00\\
Questions done & 9 & 18 & 27 & 36\\
\bottomrule
\end{tabular}
\end{center}
If you are more than two minutes behind at a checkpoint, shorten the next read to a skim of first and last sentences plus pivots, and lean harder on the question references.

\subsection{Protocol}
\begin{enumerate}
  \item \textbf{Choose your order.} Glance at all four units. If one passage type consistently costs you time (find it in your error log), save it for last.
  \item \textbf{Never leave a blank.} No penalty for guessing. In the final minute, fill every remaining question.
  \item \textbf{Flag, don't fight.} If a question takes more than 90 seconds, choose your best answer, flag it, and move on.
  \item \textbf{Treat every question as scored.} Field-test items are invisible.
  \item \textbf{Mind your stamina.} Reading comes after English and Math. Practice full sections, not just single passages, in the weeks before the test.
  \item \textbf{On the digital test,} use the highlighter for paragraph notes and the eliminator to cross out choices. Practice with the official tools in advance.
\end{enumerate}

% =====================================================================
\section{Quick Reference}
% =====================================================================
\begin{notebox}
\small
\textbf{Format.} 36 Q / 40 min / 4 units $\times$ 9 Q / $\approx$67 sec per Q / 27 scored / no guessing penalty.\\
\textbf{Loop.} Classify $\rightarrow$ go back $\rightarrow$ predict $\rightarrow$ eliminate. If you can't point to it, it's not the answer.

\vspace{4pt}
\begin{tabular}{@{}p{0.08\linewidth}p{0.85\linewidth}@{}}
\multicolumn{2}{@{}l}{\textbf{Key Ideas \& Details (44--52\%)}}\\
R1 & Details: the answer is in the text, paraphrased.\\
R2 & Inference: the smallest step the text forces you to take.\\
R3 & Main idea: covers the whole---no less, no more.\\
R4 & Relationships: map the arrow; watch for reversals.\\
\multicolumn{2}{@{}l}{\textbf{Craft \& Structure (26--33\%)}}\\
R5 & Word in context: cover it, predict, match; the common meaning is the trap.\\
R6 & Word choice and tone: what does it do that a neutral word wouldn't? Moderate tones.\\
R7 & Structure: ``why is this here?'' Name the job.\\
R8 & Purpose and POV: who, when, how they feel. Rarely neutral, rarely extreme.\\
\multicolumn{2}{@{}l}{\textbf{Integration of Knowledge \& Ideas (19--26\%)}}\\
R9 & Arguments: claim first; evidence must link directly; weaken with a confound.\\
R10 & Paired: each passage alone, then name the relationship.\\
R11 & VQI: title, axes, units, groups; level vs.\ change; shown vs.\ consistent.\\
\end{tabular}

\vspace{4pt}
\textbf{Ten traps.} Extreme $\cdot$ out of scope $\cdot$ half-right $\cdot$ misplaced detail $\cdot$ reversal $\cdot$ true but irrelevant $\cdot$ literal meaning $\cdot$ too narrow/broad $\cdot$ wrong speaker $\cdot$ consistent-not-shown.\\
\textbf{Clock.} 10 / 20 / 30 / 40. Flag at 90 seconds. Fill every blank in the last minute.
\end{notebox}

% =====================================================================
\section{Error Log}
% =====================================================================
Log every miss---from drills, from Volume~2, and from official practice tests. After 20 entries, the Rule and Trap columns will show you exactly what to study.

\begin{center}
\small
\renewcommand{\arraystretch}{1.9}
\begin{tabular}{@{}|p{0.09\linewidth}|p{0.07\linewidth}|p{0.12\linewidth}|p{0.06\linewidth}|p{0.14\linewidth}|p{0.19\linewidth}|p{0.14\linewidth}|@{}}
\hline
\textbf{Source} & \textbf{Q\#} & \textbf{Passage} & \textbf{Rule} & \textbf{Trap} & \textbf{Why I missed it} & \textbf{Fix}\\
\hline
 & & & & & & \\ \hline
 & & & & & & \\ \hline
 & & & & & & \\ \hline
 & & & & & & \\ \hline
 & & & & & & \\ \hline
 & & & & & & \\ \hline
 & & & & & & \\ \hline
 & & & & & & \\ \hline
 & & & & & & \\ \hline
 & & & & & & \\ \hline
 & & & & & & \\ \hline
 & & & & & & \\ \hline
\end{tabular}
\end{center}

\section*{Sources}
\addcontentsline{toc}{section}{Sources}
{\small
Format, reporting categories, and legacy-versus-enhanced figures: ACT, Inc., \emph{ACT Reporting Category Tables: Legacy vs.\ Enhanced} (\url{https://site.act.org}); Test Innovators, ``What To Expect on the Enhanced ACT Reading Test (2026)''; Kaplan, ``ACT Updates: Latest Changes to the Exam''; Magoosh, ``How to Get a Perfect 36 on ACT Reading (Enhanced ACT)''; Test Ninjas, ``ACT Question Types by Section.'' All drill passages are original; towns, libraries, cities, and data in them are fictional. Drill 7 paraphrases a well-known framing experiment from behavioral economics.\par}

% =====================================================================
\answerkey{Scoring: one point per correct answer. Tally by rule in your error log. Category totals in this volume: KID \arabic{catKID} $\cdot$ C\&S \arabic{catCS} $\cdot$ IKI \arabic{catIKI}.}

\end{document}
